% 海王星（综述）
% license CCBYSA3
% type Wiki

本文根据 CC-BY-SA 协议转载翻译自维基百科\href{https://en.wikipedia.org/wiki/Neptune}{相关文章}。

\begin{figure}[ht]
\centering
\includegraphics[width=8cm]{./figures/bfb6ab5163cda4f5.png}
\caption{1989年8月由旅行者2号拍摄的真实色彩海王星图像；\(^\text{[1]}\) 中央为大黑斑\(^\text{[a]}\)。} \label{fig_Neptun_1}
\end{figure}
海王星是绕太阳运行的第八颗也是已知最远的行星。按直径计算，它是太阳系中第四大行星，按质量计算为第三大行星，并且是密度最高的巨行星。它的质量是地球的17倍。与其邻近的冰巨行星天王星相比，海王星略小，但质量更大、密度更高。由于主要由气体和液体组成\(^\text{[21]}\)，它没有明确定义的固体表面。海王星每164.8年绕太阳运行一周，轨道距离为30.1个天文单位（45亿千米；28亿英里）。它以罗马海神命名，并拥有天文学符号♆，代表海神的三叉戟\(^\text{[e]}\)。

海王星肉眼不可见，也是太阳系中唯一一颗最初不是通过直接经验观测而发现的行星。相反，天王星轨道的异常变化使亚历克西·布瓦尔提出假说，认为其轨道受未知行星的引力扰动。在布瓦尔去世后，约翰·库奇·亚当斯和乌尔班·勒维耶分别依据布瓦尔的观测独立地数学预测了海王星的位置。随后在1846年9月23日\(^\text{[2]}\)，约翰·戈特弗里德·加勒用望远镜直接观测到了海王星，位置与勒维耶预测的位置相差不到一度。其最大卫星海卫一不久后被发现，但海王星剩余的其他卫星直到20世纪才通过望远镜观测到。

海王星与地球的距离使其视直径很小，而其距离太阳遥远则使其极为暗淡，给基于地面的望远镜研究带来挑战。直到哈勃太空望远镜以及配备自适应光学的大型地面望远镜出现，才使得详细观测成为可能。1989年8月25日飞掠海王星的旅行者2号仍然是唯一访问过这颗行星的航天器\(^\text{[22][23]}\)。与气体巨行星（木星和土星）类似，海王星的大气主要由氢和氦组成，并含有少量的碳氢化合物以及可能存在的氮，但其含有更高比例的冰，如水、氨和甲烷。与天王星类似，其内部主要由冰和岩石构成\(^\text{[24]}\)，这两颗行星通常被视为“冰巨行星”以示区分\(^\text{[25]}\)。除了瑞利散射外，外层区域中微量的甲烷使海王星呈现出淡蓝色\(^\text{[26][27]}\)。

与天王星强烈的季节性大气不同，后者可能长时间几乎没有可见特征，海王星的大气具有活跃且持续可见的天气形态。在旅行者2号1989年飞掠期间，该行星的南半球存在一个可与木星大红斑相比的大黑斑。2018年，新的主要暗斑及更小的暗斑被识别并研究\(^\text{[28]}\)。这些天气形态由太阳系所有行星中最强的持续风驱动，风速可高达 2{,}100 km/h（580 m/s；1{,}300 mph）\(^\text{[29]}\)。由于其与太阳的巨大距离，海王星外层大气是太阳系中最寒冷的地方之一，其云顶温度接近 55 K（−218 °C；−361 °F）。行星中心温度约为 5{,}400 K（5{,}100 °C；9{,}300 °F）\(^\text{[30][31]}\)。海王星具有微弱且破碎的环系统（称为“弧”），于 1984 年被发现并由旅行者2号确认\(^\text{[32]}\)。
\subsection{历史}
\subsubsection{发现}
\begin{figure}[ht]
\centering
\includegraphics[width=6cm]{./figures/b8d92d1f5c26ae9a.png}
\caption{} \label{fig_Neptun_2}
\end{figure}
\begin{figure}[ht]
\centering
\includegraphics[width=6cm]{./figures/a7ed05f7bd7b8f95.png}
\caption{} \label{fig_Neptun_3}
\end{figure}
\begin{figure}[ht]
\centering
\includegraphics[width=6cm]{./figures/bb642f227f35d3c1.png}
\caption{} \label{fig_Neptun_4}
\end{figure}
\begin{figure}[ht]
\centering
\includegraphics[width=6cm]{./figures/6e8babe86270c7ba.png}
\caption{} \label{fig_Neptun_5}
\end{figure}
一些已知最早的望远镜观测记录是伽利略于1612年12月28日和1613年1月27日（新历）绘制的图，它们包含的定位点与现今已知的海王星在这些日期的位置相吻合。两次观测中，伽利略似乎都把海王星误认为一颗恒星，因为它在夜空中出现在靠近木星的位置——即形成合相\(^\text{[33]}\)。因此，他不被视为海王星的发现者。在1612年12月的首次观测中，海王星在天空中几乎保持静止，因为它就在那一天开始逆行。这种视逆行运动是当地球轨道超过外行星时产生的。由于海王星只是刚刚开始其每年的逆行周期，这颗行星的运动非常微弱，以至于伽利略的小型望远镜无法检测到\(^\text{[34]}\)。2009年的一项研究表明，伽利略至少意识到他所观测到的这颗“恒星”相对于固定恒星发生了移动\(^\text{[35]}\)。

1821年，亚历克西·布瓦尔发表了关于天王星轨道的天文表\(^\text{[36]}\)。随后观测揭示出与天文表相比存在显著偏差，这促使布瓦尔假设有未知天体通过引力相互作用扰动了天王星的轨道\(^\text{[37]}\)。1843年，约翰·库奇·亚当斯开始利用他所拥有的数据研究天王星的轨道。他向皇家天文学家乔治·艾里爵士请求额外数据，艾里于1844年2月提供了这些数据。亚当斯在1845—1846年继续工作，并提出了多个关于天王星之外未发现行星位置的不同估算\(^\text{[38][39]}\)。

与亚当斯独立地，乌尔班·勒维耶在1845—1846年自行得出了指向未发现行星的计算结果，但并未激起其同胞的热情。1846年6月，在看到勒维耶首次发表的可疑未发现行星经度估算及其与亚当斯估算的相似性后，艾里说服詹姆斯·查利斯进行搜索。查利斯在整个8月和9月徒劳地搜索天空\(^\text{[37][40]}\)。事实上，查利斯比后来真正发现该行星的约翰·戈特弗里德·加勒早一年观测到了海王星，并在1845年8月4日和12日两次观测到。然而，他过时的星图和糟糕的观测技术使他在后续分析之前未能识别这些观测结果。查利斯充满悔意，但将自己的疏忽归咎于星图以及同时进行彗星观测工作的干扰\(^\text{[41][37][42]}\)。




















